\documentclass[aspectratio=169]{beamer}

\usetheme{Madrid}
\usecolortheme{whale}

\usepackage{multimedia}
\usepackage{amsmath,amssymb}
\usepackage{tikz}

% Enable notes
\setbeameroption{show notes}

\title[Speaker Tracking \& Beamer Presenter]{Robust Acoustic Speaker Tracking \\ \& Multimedia Presentation Systems}
\subtitle{Live Benchmark Deck with Beamer Notes, Video, and Audio}
\author[V. Kalkhorani et al.]{Vahid Kalkhorani\inst{1} \and Advanced Systems Group\inst{2}}
\institute[OSU]{
  \inst{1}Department of Computer Science \& Engineering, The Ohio State University \\
  \inst{2}Acoustic Signal Processing \& AI Laboratory
}
\date{\today}

\begin{document}

% -----------------------------------------------------------------------------
% SLIDE 1: Title
% -----------------------------------------------------------------------------
\begin{frame}
  \titlepage
  \note{
    \textbf{PRESENTER NOTES (Slide 1)}:
    \begin{itemize}
      \item Welcome everyone to the presentation.
      \item Introduce the objective: evaluating the next-generation Rust-based Beamer Presenter.
      \item Point out that this presentation includes embedded audio and video media.
      \item Press `L` on your keyboard to activate the laser pointer at any time.
    \end{itemize}
  }
\end{frame}

% -----------------------------------------------------------------------------
% SLIDE 2: Mathematical Formulation (Acoustic State-Space)
% -----------------------------------------------------------------------------
\begin{frame}{Mathematical Modeling: Acoustic State Space}
  \begin{block}{State-Space Equation for Moving Acoustic Source}
    Let $\mathbf{x}_k = [p_x, p_y, \dot{p}_x, \dot{p}_y]^T$ denote the state vector of the active speaker:
    \begin{equation}
      \mathbf{x}_{k} = \mathbf{F}\mathbf{x}_{k-1} + \mathbf{w}_k, \quad \mathbf{w}_k \sim \mathcal{N}(\mathbf{0}, \mathbf{Q})
    \end{equation}
  \end{block}

  \begin{exampleblock}{Microphone Array Measurement Model (TDoA)}
    For microphone pair $(i, j)$, the observed Time-Difference-of-Arrival (TDoA) is:
    \begin{equation}
      \tau_{ij}(\mathbf{x}_k) = \frac{\|\mathbf{p}_k - \mathbf{m}_i\| - \|\mathbf{p}_k - \mathbf{m}_j\|}{c} + v_{ij}
    \end{equation}
    where $c = 343\text{ m/s}$ is the acoustic sound velocity in air.
  \end{exampleblock}

  \note{
    \textbf{SPEAKER NOTES (Slide 2)}:
    \begin{itemize}
      \item Walk the audience through the continuous-to-discrete transition matrix $\mathbf{F}$.
      \item Highlight the non-linear hyperbolic nature of the TDoA measurements.
      \item If needed, use the **Pen tool (P)** to sketch the microphone array baseline.
      \item Or press `+` or `N` to open a blank whiteboard for step-by-step Taylor expansion.
    \end{itemize}
  }
\end{frame}

% -----------------------------------------------------------------------------
% SLIDE 3: Audio Media Test
% -----------------------------------------------------------------------------
\begin{frame}{Acoustic Benchmark: Audio Stimulus}
  \begin{alertblock}{Audio Playback Verification}
    Click the trigger box below to play the calibrated acoustic reference tone:
  \end{alertblock}

  \vspace{0.8em}
  \centering
  \movie[autostart=false,showcontrols=true]{
    \colorbox{blue!15}{\parbox{0.75\textwidth}{\centering \Large \vspace{0.5em} \textbf{[AUDIO] Play 440 Hz Reference Audio} \\ \small (Embedded via \texttt{\textbackslash movie} from \texttt{multimedia.sty}) \vspace{0.5em}}}
  }{sample_audio.mp3}

  \vspace{1.2em}
  \begin{itemize}
    \item Frequency: $f_0 = 440.0\text{ Hz}$ (Standard Concert Pitch $A_4$)
    \item Sample Rate: $44.1\text{ kHz}$, 16-bit PCM Audio
    \item Audio clip duration: 4.0 seconds
  \end{itemize}

  \note{
    \textbf{SPEAKER NOTES (Slide 3)}:
    \begin{itemize}
      \item Play the 440 Hz test tone to demonstrate native audio subsystem integration.
      \item Confirm that audio plays through the main presentation audio output without lag.
      \item Explain that speaker tracking algorithms initialize spectral baselines with this reference.
    \end{itemize}
  }
\end{frame}

% -----------------------------------------------------------------------------
% SLIDE 4: Video Media Test
% -----------------------------------------------------------------------------
\begin{frame}{Target Tracking: Video Stream Verification}
  \begin{block}{Live Video Tracking Simulation}
    Real-time acoustic-visual sensor fusion evaluation:
  \end{block}

  \vspace{0.4em}
  \centering
  \movie[width=8.5cm,height=4.78cm,showcontrols=true,autostart=false]{
    \colorbox{black}{\parbox{8.5cm}{\centering \textcolor{white}{\vspace{1.5cm} \textbf{[VIDEO] Click to Play Video Simulation} \\ \small sample\_video.mp4 (720p 30fps) \vspace{1.5cm}}}}
  }{sample_video.mp4}

  \note{
    \textbf{SPEAKER NOTES (Slide 4)}:
    \begin{itemize}
      \item Start the video playback.
      \item Note the zero-tearing and 30fps smooth hardware-accelerated decode.
      \item Highlight the test color patterns simulating camera sensor input.
      \item Mention that the presenter scrub bar allows seeking to any timecode seamlessly.
    \end{itemize}
  }
\end{frame}

% -----------------------------------------------------------------------------
% SLIDE 5: Conclusion & Inking Demo
% -----------------------------------------------------------------------------
\begin{frame}{Summary \& Live Inking Demonstration}
  \begin{columns}[T]
    \begin{column}{0.5\textwidth}
      \textbf{Key System Capabilities:}
      \begin{enumerate}
        \item Sub-10ms inking latency with Catmull-Rom smoothing.
        \item Instant Beamer note extraction.
        \item Multi-window display detachment.
        \item Vector-grade 600 DPI rendering quality.
      \end{enumerate}
    \end{column}
    \begin{column}{0.5\textwidth}
      \textbf{Presenter Inking Verification:}
      \begin{itemize}
        \item Use **Highlighter (H)** on these bullet points.
        \item Use **Laser (L)** to trace slide elements.
        \item Export annotations with **Save PDF**.
      \end{itemize}
    \end{column}
  \end{columns}

  \vspace{1.5em}
  \centering
  \textbf{Thank you for your attention! Questions?}

  \note{
    \textbf{SPEAKER NOTES (Slide 5)}:
    \begin{itemize}
      \item Open the floor for questions from committee or conference audience.
      \item Show how easy it is to draw derivations live on the slides or on whiteboard pages.
      \item Export the final annotated presentation to PDF when the session concludes.
    \end{itemize}
  }
\end{frame}

\end{document}
